% CP-VI-0003
% target_chapter: VI/07 — Spectrum of a Ring
% source_unit_id: IMP-DL-0E3E7EC766-R000712-68047675
% source: Downloads/theory-of-algebraic-geometry-1.tex L712-L735
% solution_provenance: SOURCE_EXPLICIT_SOLUTION

\subsection*{Problem 3}
	
	Let \(A\) be a ring. Show that the following are equivalent:
	
	\begin{enumerate}[label=\roman*)]
		\item \(\operatorname{Spec} A\) is disconnected.
		
		\item There exist nonzero elements \(e_1,e_2\in A\) such that
		\[
		e_1e_2=0,
		\qquad
		e_1^2=e_1,
		\qquad
		e_2^2=e_2,
		\qquad
		e_1+e_2=1.
		\]
		
		\item \(A\) is isomorphic to \(A_1\times A_2\) for some nonzero rings
		\(A_1,A_2\).
	\end{enumerate}
	
	\bigskip
	
	\textbf{Solution.}
	
	The idea is that disconnectedness of \(\operatorname{Spec} A\) is the same thing as the
	existence of a nontrivial idempotent decomposition of the ring.
	
	An element \(e\in A\) is called idempotent if
	\[
	e^2=e.
	\]
	The trivial idempotents are \(0\) and \(1\). A nontrivial idempotent is an
	idempotent different from \(0\) and \(1\).
	
	We prove the implications one by one.
	
	\subsection*{\((ii)\Rightarrow(i)\)}
	
	Assume that there exist nonzero elements \(e_1,e_2\in A\) such that
	\[
	e_1e_2=0,
	\qquad
	e_1^2=e_1,
	\qquad
	e_2^2=e_2,
	\qquad
	e_1+e_2=1.
	\]
	
	Consider the two basic open sets
	\[
	D(e_1)
	\qquad\text{and}\qquad
	D(e_2).
	\]
	
	Their intersection is
	\[
	D(e_1)\cap D(e_2)=D(e_1e_2).
	\]
	Since \(e_1e_2=0\), this becomes
	\[
	D(e_1)\cap D(e_2)=D(0).
	\]
	But
	\[
	D(0)=\{\mathfrak p\in\operatorname{Spec} A\mid 0\notin \mathfrak p\}=\varnothing,
	\]
	because every ideal contains \(0\). Hence
	\[
	D(e_1)\cap D(e_2)=\varnothing.
	\]
	
	Now we show that
	\[
	D(e_1)\cup D(e_2)=\operatorname{Spec} A.
	\]
	Let \(\mathfrak p\in\operatorname{Spec} A\). Suppose, for contradiction, that
	\[
	\mathfrak p\notin D(e_1)\cup D(e_2).
	\]
	Then
	\[
	e_1\in\mathfrak p
	\quad\text{and}\quad
	e_2\in\mathfrak p.
	\]
	Therefore
	\[
	e_1+e_2\in \mathfrak p.
	\]
	But
	\[
	e_1+e_2=1.
	\]
	Thus
	\[
	1\in\mathfrak p,
	\]
	which is impossible because a prime ideal is proper. Hence
	\[
	D(e_1)\cup D(e_2)=\operatorname{Spec} A.
	\]
	
	It remains to see that both \(D(e_1)\) and \(D(e_2)\) are nonempty.
	
	If \(D(e_1)=\varnothing\), then by Problem 1, \(e_1\) is nilpotent. But \(e_1\)
	is idempotent:
	\[
	e_1^2=e_1.
	\]
	If an idempotent is nilpotent, then it must be zero. Indeed, if \(e_1^n=0\), then
	for every \(n\geq 1\),
	\[
	e_1^n=e_1.
	\]
	Thus \(e_1=0\), contradicting the assumption that \(e_1\neq 0\). Therefore
	\[
	D(e_1)\neq \varnothing.
	\]
	Similarly,
	\[
	D(e_2)\neq \varnothing.
	\]
	
	Therefore
	\[
	\operatorname{Spec} A=D(e_1)\sqcup D(e_2)
	\]
	is a decomposition of \(\operatorname{Spec} A\) into two disjoint nonempty open subsets. Hence
	\(\operatorname{Spec} A\) is disconnected.
	
	\subsection*{\((ii)\Rightarrow(iii)\)}
	
	Assume again that
	\[
	e_1e_2=0,
	\qquad
	e_1^2=e_1,
	\qquad
	e_2^2=e_2,
	\qquad
	e_1+e_2=1.
	\]
	
	Define
	\[
	A_1=Ae_1,
	\qquad
	A_2=Ae_2.
	\]
	These are ideals of \(A\), but they are also rings in their own right. The
	identity element of \(Ae_1\) is \(e_1\), and the identity element of \(Ae_2\) is
	\(e_2\).
	
	Define a map
	\[
	\varphi:A\longrightarrow Ae_1\times Ae_2
	\]
	by
	\[
	\varphi(a)=(ae_1,ae_2).
	\]
	
	We show that \(\varphi\) is a ring isomorphism.
	
	First, \(\varphi\) is additive:
	\[
	\varphi(a+b)=((a+b)e_1,(a+b)e_2)
	=(ae_1+be_1,ae_2+be_2)
	=\varphi(a)+\varphi(b).
	\]
	
	It is multiplicative:
	\[
	\varphi(ab)=(abe_1,abe_2).
	\]
	On the other hand,
	\[
	\varphi(a)\varphi(b)
	=
	(ae_1,ae_2)(be_1,be_2)
	=
	(ae_1be_1,ae_2be_2).
	\]
	Since \(e_1^2=e_1\) and \(e_2^2=e_2\), we get
	\[
	ae_1be_1=ab e_1^2=ab e_1,
	\]
	and
	\[
	ae_2be_2=ab e_2^2=ab e_2.
	\]
	Therefore
	\[
	\varphi(a)\varphi(b)=\varphi(ab).
	\]
	
	Also,
	\[
	\varphi(1)=(e_1,e_2),
	\]
	which is the identity element of \(Ae_1\times Ae_2\).
	
	Now define
	\[
	\psi:Ae_1\times Ae_2\longrightarrow A
	\]
	by
	\[
	\psi(x,y)=x+y.
	\]
	
	Let \(a\in A\). Then
	\[
	\psi(\varphi(a))
	=
	\psi(ae_1,ae_2)
	=
	ae_1+ae_2
	=
	a(e_1+e_2)
	=
	a.
	\]
	Thus
	\[
	\psi\circ\varphi=\operatorname{id}_A.
	\]
	
	Conversely, let \(x\in Ae_1\) and \(y\in Ae_2\). Then there exist \(a,b\in A\)
	such that
	\[
	x=ae_1,
	\qquad
	y=be_2.
	\]
	We have
	\[
	xe_1=x,
	\qquad
	xe_2=0,
	\]
	and
	\[
	ye_1=0,
	\qquad
	ye_2=y.
	\]
	Therefore
	\[
	\varphi(\psi(x,y))
	=
	\varphi(x+y)
	=
	((x+y)e_1,(x+y)e_2)
	=
	(xe_1+ye_1,xe_2+ye_2)
	=
	(x,y).
	\]
	Thus
	\[
	\varphi\circ\psi=\operatorname{id}_{Ae_1\times Ae_2}.
	\]
	
	Hence \(\varphi\) is an isomorphism, and
	\[
	A\cong Ae_1\times Ae_2.
	\]
	Since \(e_1,e_2\neq 0\), both rings \(Ae_1\) and \(Ae_2\) are nonzero.
	
	Therefore condition \((iii)\) holds.
	
	\subsection*{\((iii)\Rightarrow(ii)\)}
	
	Assume that
	\[
	A\cong A_1\times A_2,
	\]
	where \(A_1\) and \(A_2\) are nonzero rings.
	
	Inside \(A_1\times A_2\), define
	\[
	e_1=(1,0),
	\qquad
	e_2=(0,1).
	\]
	Then
	\[
	e_1^2=(1,0)^2=(1,0)=e_1,
	\]
	and
	\[
	e_2^2=(0,1)^2=(0,1)=e_2.
	\]
	Also,
	\[
	e_1e_2=(1,0)(0,1)=(0,0),
	\]
	and
	\[
	e_1+e_2=(1,0)+(0,1)=(1,1),
	\]
	which is the identity element of \(A_1\times A_2\).
	
	Because \(A_1\) and \(A_2\) are nonzero, we have
	\[
	e_1\neq 0,
	\qquad
	e_2\neq 0.
	\]
	Transporting these elements through the isomorphism
	\[
	A\cong A_1\times A_2,
	\]
	we obtain nonzero elements \(e_1,e_2\in A\) satisfying
	\[
	e_1e_2=0,
	\qquad
	e_1^2=e_1,
	\qquad
	e_2^2=e_2,
	\qquad
	e_1+e_2=1.
	\]
	Therefore condition \((ii)\) holds.
	
	\subsection*{\((i)\Rightarrow(ii)\)}
	
	Assume that \(\operatorname{Spec} A\) is disconnected. Then there exist nonempty disjoint open
	subsets \(U,V\subseteq \operatorname{Spec} A\) such that
	\[
	\operatorname{Spec} A=U\sqcup V.
	\]
	Since \(V=\operatorname{Spec} A\setminus U\), the set \(U\) is also closed. Similarly, \(V\) is
	also closed. Hence \(U\) and \(V\) are both open and closed, that is, clopen.
	
	We use the following standard fact.
	
	\medskip
	
	\noindent
	\textbf{Lemma.}
	For a ring \(A\), clopen subsets of \(\operatorname{Spec} A\) correspond exactly to
	idempotents of \(A\). More precisely, if \(e\in A\) is idempotent, then
	\[
	D(e)
	\]
	is clopen. Conversely, every clopen subset of \(\operatorname{Spec} A\) is of the form
	\[
	D(e)
	\]
	for some idempotent \(e\in A\).
	
	\medskip
	
	Let us first check the easy direction of the lemma. Suppose \(e^2=e\). Then
	\[
	e(1-e)=0.
	\]
	For a prime ideal \(\mathfrak p\), since
	\[
	e(1-e)\in\mathfrak p,
	\]
	we must have
	\[
	e\in\mathfrak p
	\quad\text{or}\quad
	1-e\in\mathfrak p.
	\]
	Moreover, both cannot occur, because then
	\[
	1=e+(1-e)\in\mathfrak p,
	\]
	which is impossible.
	
	Thus, for every prime ideal \(\mathfrak p\), exactly one of \(e\) and \(1-e\)
	belongs to \(\mathfrak p\). Hence
	\[
	D(e)=V(1-e).
	\]
	Indeed,
	\[
	\mathfrak p\in D(e)
	\]
	means
	\[
	e\notin\mathfrak p,
	\]
	which is equivalent to
	\[
	1-e\in\mathfrak p.
	\]
	Therefore
	\[
	D(e)=V(1-e).
	\]
	Thus \(D(e)\) is open by definition and closed because it is equal to
	\(V(1-e)\). Hence \(D(e)\) is clopen.
	
	Conversely, every clopen subset arises in this way from an idempotent. This is
	the standard idempotent--clopen correspondence for affine schemes.
	
	Since \(U\) is clopen, there exists an idempotent \(e_1\in A\) such that
	\[
	U=D(e_1).
	\]
	Define
	\[
	e_2=1-e_1.
	\]
	Then
	\[
	e_2^2=(1-e_1)^2
	=
	1-2e_1+e_1^2.
	\]
	Since \(e_1^2=e_1\), this becomes
	\[
	e_2^2
	=
	1-2e_1+e_1
	=
	1-e_1
	=
	e_2.
	\]
	Therefore \(e_2\) is also idempotent.
	
	Moreover,
	\[
	e_1e_2=e_1(1-e_1)=e_1-e_1^2=0,
	\]
	and
	\[
	e_1+e_2=e_1+(1-e_1)=1.
	\]
	
	It remains to check that \(e_1\) and \(e_2\) are nonzero.
	
	Because \(U\neq\varnothing\), we have
	\[
	D(e_1)\neq\varnothing.
	\]
	If \(e_1=0\), then
	\[
	D(e_1)=D(0)=\varnothing,
	\]
	contradicting \(U\neq\varnothing\). Hence
	\[
	e_1\neq 0.
	\]
	
	Similarly, the complement of \(D(e_1)\) is
	\[
	V(e_1).
	\]
	Since \(e_1\) is idempotent, this complement is also
	\[
	D(1-e_1)=D(e_2).
	\]
	But this complement is \(V\), and \(V\neq\varnothing\). Therefore
	\[
	D(e_2)\neq\varnothing.
	\]
	Hence
	\[
	e_2\neq 0.
	\]
	
	Thus we have found nonzero elements \(e_1,e_2\in A\) such that
	\[
	e_1e_2=0,
	\qquad
	e_1^2=e_1,
	\qquad
	e_2^2=e_2,
	\qquad
	e_1+e_2=1.
	\]
	Therefore condition \((ii)\) holds.
	
	\subsection*{Conclusion}
	
	We have proved
	\[
	(ii)\Rightarrow(i),
	\qquad
	(ii)\Rightarrow(iii),
	\qquad
	(iii)\Rightarrow(ii),
	\qquad
	(i)\Rightarrow(ii).
	\]
	Therefore the three conditions are equivalent:
	\[
	\boxed{
		\operatorname{Spec} A \text{ is disconnected}
		\iff
		A \text{ has a nontrivial idempotent decomposition}
		\iff
		A\cong A_1\times A_2
	}
	\]
	for some nonzero rings \(A_1,A_2\).
